\subsection{The residue-characteristic check}

\begin{result}{E.X1}{Refuted; precise comparison proved}{Full completion at residue characteristics two and three}
For $p=2$ and $p=3$, put
$U_p=\mathbf P^1_{\overline{\F}_p}-\{0,1,\infty\}$.
The full geometric etale fundamental group of $U_p$ is
not isomorphic to the full completion of the topological
tripod group $F_2$.  More precisely, the explicit covers
\[
                  y^p-y=t^m,\qquad m>0,\quad p\nmid m,
 \tag{E.X1}
\]
give infinitely many linearly independent continuous
$C_p$-characters of $\pi_1^{\mathrm{et}}(U_p)$.
The exact comparison that survives is
\[
 \widehat F_2^{(p')}\xrightarrow{\ \sim\ }
             \pi_1^{\mathrm{et}}(U_p)^{(p')},
 \tag{E.X2}
\]
where the superscript denotes the maximal pro-prime-to-$p$
quotient.  It comes from specialization of the relative
pair $(\mathbf P^1,\{0,1,\infty\})$ over the Witt ring,
is defined up to inner automorphism after basepoints and
a specialization path are chosen, and carries each of
the three inertia groups to its corresponding inertia
group.  Thus the full-completion statements over
$\overline\Q$ used in this paper have no unchanged
specialization to either of these residue characteristics.
\end{result}

\begin{proof}
Set $k=\overline{\F}_p$ and
$R=k[t,t^{-1},(t-1)^{-1}]$.
The algebra $R[y]/(y^p-y-t^m)$ is finite free of rank
$p$ and etale, since its derivative with respect to
$y$ is $-1$.  Translation $y\mapsto y+1$ gives its
$C_p$ action.  The Artin--Schreier exact sequence
identifies these torsor classes with $[t^m]$ in
$R/\{f^p-f:f\in R\}$; this exact sequence and its
cohomological proof were checked in
\cite[Tag 0A3J]{StacksArtinSchreier} and the explicit
polynomial description in
\cite[XIII, p.~292]{SGA1Revised}.

Consider a nonzero finite $\F_p$-linear combination of
the indicated monomials and let $M$ be its largest
exponent.  Its only pole is at infinity and has order
$M$, which is prime to $p$.  An equality with $f^p-f$
would first force $f$ to have no poles at zero or one:
a pole of order $a$ there gives a pole of order $pa$
in $f^p-f$, with no cancellation of the leading term.
Thus $f$ is a polynomial.  A nonconstant such polynomial
gives a pole at infinity of order $p\deg f$, again
with no cancellation; a constant gives no pole.
Both contradict the pole of order $M$.  This proves
the asserted $\F_p$-linear independence in the quotient
formed over $k$, including all possible $f\in R$.
In particular each nonzero class defines
a connected cover.  The rank-two free profinite group
has exactly a two-dimensional space of continuous
$C_p$-characters, since a character is determined by
the two free generators.  An isomorphism with the
displayed geometric group would induce an isomorphism
of these character spaces, which the infinite family
rules out.

For (E.X2) take $W(k)$ and the smooth proper relative
curve $\mathbf P^1_{W(k)}$ with its three disjoint
sections.  The hypotheses of SGA~1 XIII, Corollary
2.9 and the specialization construction 2.10 are
therefore satisfied for finite groups of order prime
to $p$.  The proof of Corollary~2.12, read in full,
identifies the induced specialization on these
quotients as an isomorphism; Lemma~2.11 identifies
the images of the three inertia groups.  On the
geometric characteristic-zero fibre, Riemann existence
and characteristic-zero base extension give
$\widehat F_2$.  Equivalently, Corollary~2.12 gives
three inertia generators with the sole relation
$c_0c_1c_\infty=1$ in the pro-prime-to-$p$ category.
Eliminating $c_\infty$ proves exactly (E.X2).
All choices and the resulting inner ambiguity are
those in the source's construction 2.10; the full
Artin--Schreier covers above are outside this
prime-to-$p$ quotient.

The program \path{code/e_residue_checks.py} performs
finite-field row reduction on the monomials and the
images $t^{pi}-t^i$ through degrees $6,12,24,48$,
separately for $p=2,3$.  Its saved output checks the
independence and the nonzero etale derivative in
both requested characteristics.  The pole argument
proves the unlimited-degree assertion.  The use of
short Weierstrass forms and of $\mu_2$-band
calculations in E.T4--E.T7 is over $\overline\Q$;
the explicit morphism (E.X2) is the residue comparison
asserted here.
\end{proof}
